% table-style.tex --- shared cell-tint vocabulary for the Experiments
% chapter's tables, defined once and used by every section (the legend
% lives in the conventions section, sec:exp-conventions).
%
% Requires: xcolor + colortbl in the including document's preamble
% (chapter-standalone.tex and the book's preamble both load them).
% The accent palette is the book's; \providecolor keeps these
% definitions inert wherever the book has already defined them.
\providecolor{accent}{HTML}{F1D9CB}
\providecolor{accentmid}{HTML}{C4623F}
\providecolor{NoteGray}{RGB}{92,92,92}
\providecolor{MarkGreen}{RGB}{22,122,58}

% The win-tint ladder: how far below the comparison baseline the
% certified value sits, in orders of magnitude of the improvement
% ratio.  Color is always decoration on top of the printed ratio ---
% never the only carrier of the distinction.
\colorlet{winA}{accent!22}     % within one order of magnitude
\colorlet{winB}{accent!45}     % 1--10 orders
\colorlet{winC}{accentmid!30}  % 10--25 orders
\colorlet{winD}{accentmid!48}  % 25--45 orders
\colorlet{winE}{accentmid!68}  % 45+ orders

% exact-cell tint: a certified value with zero conservatism against
% its model (paired with the cell's own text, e.g. "optimal, exactly")
\colorlet{exacttint}{MarkGreen!14}

% legend swatch (text-mode; no TikZ dependency)
\providecommand{\swx}[1]{\begingroup\fboxsep0pt\colorbox{#1}{\rule{0pt}{1.6ex}\hspace{2.4mm}}\endgroup}

% alternate-reading row marker + demphasized table text
\providecommand{\altrow}{{\footnotesize$\hookrightarrow$}\,}
\providecommand{\tdemph}[1]{\textcolor{NoteGray}{#1}}
